\chapter{SOAP vs REST}
\label{cap:soap-vs-rest}

% TODO: integrar el aporte validado de Carvajal Loor Johan Stalin
% (investigacion comparativa SOAP vs REST, minimo 8 criterios, casos de
% uso aplicables, ejemplo SOAP y su equivalente REST).
%
% Fuente prevista (no integrada todavia en este capitulo):
%   docs/u4/investigacion/SOAP-VS-REST.md
%
% Al integrar, citar como minimo:
%   \cite{fielding2000} -- Fielding (2000), REST como estilo arquitectonico
%   \cite{w3c-soap12}   -- especificacion oficial SOAP 1.2 (W3C)
% y anadir aqui las entradas bibliograficas adicionales que Carvajal
% verifique para SOAP/REST (sin inventar autores, anios ni URLs).

\pendienteIntegracion{Investigaci\'on comparativa SOAP vs REST (Carvajal Loor Johan Stalin). Este cap\'itulo se completar\'a con el aporte validado, incluyendo al menos 8 criterios de comparaci\'on, casos de uso aplicables a cada estilo, y un ejemplo SOAP con su equivalente REST.}

\section{Alcance previsto}
\label{sec:soap-rest-alcance}

Este cap\'itulo integrar\'a, cuando est\'e disponible, la investigaci\'on
comparativa entre los estilos arquitect\'onicos SOAP y REST, con al menos
ocho criterios de comparaci\'on (por ejemplo: protocolo de transporte,
formato de mensaje, contrato de interfaz, manejo de estado, cach\'e,
seguridad, tolerancia a fallos y curva de adopci\'on), casos de uso donde
cada estilo resulta m\'as apropiado, y un ejemplo concreto de una operaci\'on
expresada en SOAP junto con su equivalente REST real de BIOPET (por
ejemplo, \texttt{GET /api/mascotas}).

No se redacta contenido definitivo de comparaci\'on en esta versi\'on del
informe porque depende de la investigaci\'on independiente asignada a
Carvajal Loor Johan Stalin, todav\'ia no incorporada a este documento.
